\documentclass[../../../include/open-logic-section]{subfiles}

\begin{document}

\olfileid{sth}{choice}{banach}
\olsection{تناقض‌نمای باناخ--تارسکی}

همچنین شاید بکوشیم انتخاب را، همان‌گونه که بولوس کوشید جایگزینی را
توجیه کند، با توسل به ملاحظات \emph{بیرونی} توجیه کنیم (نگاه کنید
به \olref[replacement][extrinsic]{sec}). آخر، پذیرفتن انتخاب پیامدهای
مطلوب بسیاری دارد: توانایی مقایسهٔ هر دو کاردینال؛ توانایی
خوش‌ترتیب‌کردن هر مجموعه؛ توانایی در نظر گرفتن کاردینال‌ها به‌عنوان
نوع خاصی از عدد ترتیبی؛ و جز این‌ها.

بااین‌حال، گاه ادعا می‌شود که انتخاب پیامدهایی \emph{نامطلوب} دارد.
این ادعا بیشتر به سبب نتیجه‌ای از \cite{BanachTarski1924} است.

\begin{thm}[تناقض‌نمای باناخ--تارسکی (در $\ZFC$)]
هر گوی را می‌توان به شمار متناهی‌ای قطعه تجزیه کرد و آن قطعه‌ها را
(با دوران و انتقال) دوباره چنان کنار هم گذاشت که دو نسخه از همان
گوی ساخته شود.
\end{thm}
\noindent
در نگاه نخست، این نتیجه تا اندازه‌ای شگفت‌انگیز است. آشکارا دو گوی
\emph{دو برابرِ} گوی اصلی حجم دارند. اما حرکت‌های صلب — دوران و
انتقال — حجم را تغییر نمی‌دهند. پس به نظر می‌رسد باناخ--تارسکی به
ما اجازه می‌دهد مادهٔ تازه‌ای را به‌گونه‌ای جادویی پدید آوریم.

اوضاع از این هم بدتر می‌شود.\footnote{نگاه کنید به
\citet[Theorem 3.12]{Wagon2016}.} استدلالی مشابه نشان می‌دهد که
می‌توان نخودی را به شمار متناهی‌ای قطعه برید و سپس قطعه‌ها را (با
دوران و انتقال) دوباره چنان کنار هم گذاشت که شیئی هم‌شکل و
هم‌اندازهٔ برج بیگ بن ساخته شود.

اما هیچ‌یک از این مطالب تنها در $\ZF$ برقرار نیست.\footnote{البته
باناخ--تارسکی را می‌توان با اصولی اکیداً ضعیف‌تر از انتخاب اثبات
کرد؛ نگاه کنید به \citet[303]{Wagon2016}.} پس با تصمیمی روبه‌رو
هستیم: انتخاب را رد کنیم، یا بیاموزیم با این «تناقض‌نما» کنار
بیاییم.

پیشنهاد ما این است که باید بیاموزیم با این «تناقض‌نما» کنار بیاییم.
در واقع، به نظر ما اصلاً تناقض‌نمای چندانی در کار نیست. به‌ویژه روشن
نیست که چرا این نتیجه باید از هریک از نتایج زیر بیشتر یا کمتر
تناقض‌نما باشد:\footnote{\citet[276--7]{Potter2004}،
\citet[16]{Weston2003} و \citet[31, 308--9]{Wagon2016} با
نمونه‌هایی دیگر نکاتی مشابه را مطرح می‌کنند.}
\begin{enumerate}
	\item شمار نقاط بازهٔ $(0,1)$ با شمار نقاط $\Real$ برابر است.
	\\\emph{اثبات}: $\tan(\pi(r-\nicefrac{1}{2}))$ را در نظر بگیرید.
	\item شمار نقاط یک خط با شمار نقاط یک مربع برابر است.
	\\نگاه کنید به \olref[his][set][pathology]{sec} و
	\olref[his][set][cantorplane]{sec}.
	\item منحنی‌های پُرکنندهٔ فضا وجود دارند.
	\\نگاه کنید به \olref[his][set][pathology]{sec} و
	\olref[his][set][hilbertcurve]{sec}.
\end{enumerate}
هیچ‌یک از این سه نتیجه به انتخاب نیاز ندارد. در واقع، اکنون آن‌ها
را صرفاً بخش‌هایی شگفت‌انگیز و دوست‌داشتنی از ریاضیات می‌دانیم.
شاید باید دربارهٔ تناقض‌نمای باناخ--تارسکی نیز همین نگرش را داشته
باشیم.

البته در اینجا یک نکتهٔ فنی لازم است؛ اما تنها چیزی که می‌طلبد حفظ
خونسردی است. حرکت‌های صلب حجم را حفظ می‌کنند. در نتیجه، هر پنج
قطعه‌ای\footnote{تناقض‌نما را بر حسب «شمار متناهی‌ای قطعه» بیان
کردیم. در واقع، \citet{Robinson1947} اثبات کرد که این تجزیه با
\emph{پنج} قطعه (و نه کمتر) شدنی است. برای اثبات، نگاه کنید به
\citet[pp.~66--7]{Wagon2016}.} که گوی به آن‌ها تجزیه می‌شود
نمی‌توانند همگی \emph{اندازه‌پذیر} باشند. به بیان تقریبی، نسبت‌دادن
حجم به هریک از این قطعه‌ها معنا ندارد. باید آن‌ها را «پراکندگی‌های
نامتناهیِ» تجسم‌ناپذیری از نقاط بدانید. شاید تصور چنین مجموعه‌های
«بی‌نهایت پراکنده‌ای» عجیب باشد. اما به نظر می‌رسد وجود آن‌ها از
فرمانی نتیجه می‌شود که در \stagesacc{} متجسد است: باید \emph{همهٔ
گردایه‌های ممکن} از مجموعه‌های پیش‌تر در دسترس را تشکیل دهید.

اگر هیچ‌یک از این مطالب قانع‌کننده نیست، استدلال (بیرونیِ) نهایی زیر
را به سود پذیرفتن تناقض‌نمای باناخ--تارسکی در نظر بگیرید. این
تناقض‌نما بی‌درنگ بهترین لطیفهٔ ریاضی تاریخ را نتیجه می‌دهد:
\begin{enumerate}
	\item[] \emph{پرسش}. آناگرام «Banach--Tarski» چیست؟
	\item[] \emph{پاسخ}. «Banach--Tarski Banach--Tarski».
\end{enumerate}

\end{document}
